\appendixitem{将紧群嵌入李群的乘积}

命题 1. 每个紧拓扑群 G 都同构于某个紧李群乘积的闭子群。

记 \(\hat{\mathbf{G}}\) 为 \(\mathbf{G}\) 在有限维复希尔伯特空间上的不可约连续酉表示的等价类集合（Spectral Theory, in preparation）。对于所有 \(u\in \hat{\mathbf{G}}\)，令 \(\mathrm{H}_u\) 为 \(u\) 的表示空间，并令 \(\rho_{u}:\mathbf{G}\to \mathbf{U}(\mathrm{H}_{u})\) 为与 \(u\) 相关的同态。由 Peter–Weyl 定理（Spectral Theory, in preparation），连续同态 \(\rho = (\rho_u)_{u\in \hat{\mathbf{G}}}\) 从 \(\mathbf{G}\) 到 \(\prod_{u\in \hat{\mathbf{G}}}\mathbf{U}(\mathrm{H}_u)\) 是单射；由于 \(\mathbf{G}\) 紧，\(\rho\) 诱导出从 \(\mathbf{G}\) 到群 \(\prod_{\hat{\mathbf{G}}} \mathbf{U}(\mathrm{H}_u)\) 的一个闭子群的同构。
推论 1. 设 V 是 G 的单位元的一个邻域。则 V 包含 G 的一个闭正规子群 H，使得商群 G/H 是李群。

设 \((\mathrm{K}_{\lambda})_{\lambda \in \mathrm{L}}\) 是一族紧李群，使得 \(\mathrm{G}\) 可与 \(\prod_{\lambda \in \mathrm{L}}\mathrm{K}_{\lambda}\) 的一个闭子群识别；对于 \(\lambda \in \mathrm{L}\)，记 \(p_{\lambda}:\mathrm{G}\to \mathrm{K}_{\lambda}\) 为典则投影在 G 上的限制。存在有限子集 \(J \subset L\)，并且对于每个 \(\lambda \in J\)，存在一个邻域 \(V_{\lambda}\)，它是 \(K_{\lambda}\) 中原点的邻域，使得 V 包含 \(\bigcap_{\lambda \in J} p_{\lambda}^{-1}(V_{\lambda})\)。现在只需置 \(H = \bigcap_{\lambda \in J} \text{Ker}(p_{\lambda})\)。

记 \((\mathrm{H}_{\alpha})_{\alpha\in\mathrm{I}}\) 为 G 的闭正规子群的递降滤族，使得商群 \(G/H_{\alpha}\) 是李群。考虑紧李群 \(G/H_{\alpha}\) 所成的射影系统（cf.~General Topology, Chap. III, §7, no. 2, Prop. 2）。

推论 2. 典则映射 \(\mathrm{G} \to \varprojlim_{\alpha} \mathrm{G} / \mathrm{H}_{\alpha}\) 是拓扑群同构。

事实上，推论 1 推出 General Topology, Chap. III, §7, no. 2 的条件 (PA) 得到满足；该断言现由 loc. cit. 的命题 2 推出。

推论 3. G 是李群，当且仅当 G 的单位元 e 存在一个邻域，该邻域不包含异于 \(\{e\}\) 的正规子群。

该条件的必要性已经证明（Chap. III, §4, no. 2, Cor. 1 of Th. 2），而充分性是推论 1 的直接结果。

